% Part: normal-modal-logic
% Chapter: frame-correspondence
% Section: introduction

\documentclass[../../../include/open-logic-section]{subfiles}

\begin{document}

\olfileid[hi]{nml}{frd}{int}

\olsection{परिचय}

मोडल तर्कशास्त्रियों की रुचि के प्रश्नों में एक है अभिगम्यता संबंध और उस
अभिगम्यता संबंध वाले प्रतिरूपों में कुछ !!{formula}ों की सत्यता के बीच का
संबंध। उदाहरणार्थ, मान लें अभिगम्यता संबंध स्वपरावर्ती है, अर्थात् प्रत्येक
$w \in W$ के लिए $Rww$। दूसरे शब्दों में, प्रत्येक जगत स्वयं से अभिगम्य है।
इसका अर्थ है कि जब $\Box !A$ किसी जगत~$w$ पर सत्य हो, तो स्वयं $w$ उन
अभिगम्य जगतों में है जिन पर $!A$ को सत्य होना चाहिए। अतः यदि
$\mModel{M}$ का अभिगम्यता संबंध~$R$ स्वपरावर्ती है, तो हम चाहे जो जगत $w$
और सूत्र $!A$ लें, $\Box !A \lif !A$ वहाँ सत्य होगा (दूसरे शब्दों में,
स्कीमा $\Box p \lif p$ और उसके सभी प्रतिस्थापन-दृष्टांत $\mModel{M}$ में
सत्य हैं)।

किंतु विलोम असत्य है। उदाहरणार्थ, ऐसा नहीं है कि यदि $\Box p \lif p$,
$\mModel{M}$ में सत्य हो, तो $R$ स्वपरावर्ती हो। हम आसानी से ऐसा अस्वपरावर्ती
प्रतिरूप~$\mModel{M}$ पा सकते हैं जिसमें $\Box p \lif p$ सभी जगतों पर सत्य
हो: केवल एक जगत~$w$ वाला प्रतिरूप लें, जिसमें $w$ स्वयं से अभिगम्य नहीं है,
पर $w \in V(p)$। $p$ का सत्य-मूल्य उपयुक्त रूप से चुनकर हम
$\Box !A \lif !A$ को ऐसे प्रतिरूप में सत्य बना सकते हैं जो स्वपरावर्ती नहीं
है।

समाधान यह है कि चर-नियतन~$V$ को समीकरण से हटा दिया जाए। यदि हम अपेक्षा करें
कि $V(p)$ में चाहे जो जगत हों, $\Box p \lif p$, $\mModel{M}$ के सभी जगतों
पर सत्य हो, तो $R$ का स्वपरावर्ती होना आवश्यक है। क्योंकि किसी भी
अस्वपरावर्ती प्रतिरूप में कम-से-कम एक जगत~$w$ ऐसा होगा कि $Rww$ नहीं। यदि हम
$V(p) = W \setminus \{w\}$ रखें, तो $p$, $w$ के अतिरिक्त सभी जगतों पर और
इसलिए $w$ से अभिगम्य सभी जगतों पर सत्य होगा (क्योंकि $w$ स्वयं से अभिगम्य
नहीं है और केवल $w$ पर $p$ असत्य है)। दूसरी ओर, $p$, $w$ पर असत्य है, अतः
$\Box p \lif p$, $w$ पर असत्य है।

इससे संकेत मिलता है कि मूल्यांकन रहित प्रतिरूप-संरचनाओं के लिए संकेत-पद्धति
प्रस्तुत करनी चाहिए; इन्हें हम \emph{फ़्रेम} कहते हैं। फ़्रेम
$\mModel{F}$ केवल युग्म $\tuple{W, R}$ है, जिसमें अभिगम्यता संबंध सहित जगतों
का समुच्चय होता है। तब प्रत्येक प्रतिरूप $\tuple{W, R, V}$ को हम फ़्रेम
$\tuple{W, R}$ पर \emph{आधारित} कहते हैं। इसके विपरीत, कोई फ़्रेम उस पर
आधारित प्रतिरूपों के वर्ग को निर्धारित करता है; और फ़्रेमों का कोई वर्ग,
उस वर्ग के किसी फ़्रेम पर आधारित प्रतिरूपों का वर्ग निर्धारित करता है। हम
$\mModel{F} \Entails !A$, किसी !!{formula} के फ़्रेम में \emph{मान्य} होने की
धारणा, इस प्रकार परिभाषित कर सकते हैं: $\mModel{F}$ पर आधारित प्रत्येक
$\mModel{M}$ के लिए $\mSat{M}{!A}$।

इस संकेत-पद्धति से हम !!{formula}ों और फ़्रेमों के वर्गों के बीच अनुरूपता संबंध
स्थापित कर सकते हैं: उदाहरणार्थ, $\mModel{F} \Entails \Box p \lif p$ तभी
जब $\mModel{F}$ स्वपरावर्ती है।

\end{document}
